% ============================================================
%  SECCIÓN: Obtención de datos y construcción del banco de pruebas
%  Formato IEEE (IEEEtran) — español
%  Complementa a benchmark_section.tex
% ============================================================

\section{Obtención de Datos}
\label{sec:datos}

El sistema distingue dos necesidades de datos con propósitos distintos:
material de \emph{entrenamiento} para el modelo generativo y un \emph{banco
de pruebas} independiente para la evaluación cuantitativa
(Sección~\ref{sec:benchmark}).
La interfaz expone ambos caminos de forma explícita, de modo que el usuario
puede \textbf{reutilizar pistas ya descargadas} o \textbf{descargar material
nuevo} sin salir de la aplicación.

\subsection{Bancos de datos disponibles}
\label{sec:datos:bancos}

\begin{itemize}
  \item \textbf{MTG-Jamendo}~\cite{bogdanov2019mtg}: catálogo de música con
        licencia Creative Commons, etiquetada por género. Es la fuente de
        \emph{entrenamiento}. En el frontend, el paso \textit{Importar audio}
        ofrece dos opciones:
        \begin{itemize}
          \item \textit{Opción A — Usar pistas ya descargadas}: un selector
                lista los catálogos locales con sus conteos por género, de
                modo que el usuario reutiliza el material previamente bajado.
          \item \textit{Opción B — Descargar nuevas pistas}: un panel permite
                indicar géneros, \textit{tags} y cantidad, y descargar audio
                nuevo desde la CDN de Jamendo.
        \end{itemize}
  \item \textbf{MusicCaps}~\cite{agostinelli2023musiclm}: 5\,521 clips de
        YouTube de 10\,s, cada uno anotado con un \textit{caption} detallado
        en inglés producido por músicos profesionales. Es la fuente del
        \emph{banco de pruebas} de evaluación. Se construye con el guion
        reproducible \texttt{scripts/sample\_testset.py}
        (Sección~\ref{sec:datos:testset}).
  \item \textbf{Audio individual}: una pista local única, para análisis
        puntual fuera del flujo de entrenamiento o evaluación.
\end{itemize}

% ----------------------------------------------------------
\subsection{Fuente de datos del banco de pruebas: MusicCaps}
\label{sec:datos:musiccaps}

MusicCaps es un subconjunto de AudioSet disponible públicamente a través de
Google Research.\footnote{\url{https://huggingface.co/datasets/google/MusicCaps}}
Cada \textit{caption} incluye género, instrumentación, tempo, estado de ánimo
y estructura, lo que lo hace idóneo para evaluar el condicionamiento textual
de modelos generativos de música.

% ----------------------------------------------------------
\subsection{Construcción del banco de pruebas}
\label{sec:datos:testset}

\subsubsection{Criterios de selección}

Se seleccionan 100 clips que cumplen los siguientes criterios:

\begin{enumerate}
  \item El video de YouTube está disponible para descarga (verificado con
        \texttt{yt-dlp -{}-simulate}).
  \item El clip se normaliza a una duración exacta de 10\,s a 32\,kHz.
  \item El \textit{caption} tiene al menos 20 palabras y menciona, como
        mínimo, género e instrumento, junto con una referencia de tempo.
  \item El conjunto es balanceado: 10 géneros $\times$ 10 clips
        (Jazz, Rock, Electrónica, Clásica, Pop, Hip-Hop, R\&B, Metal, Folk,
        Latina).
\end{enumerate}

\paragraph{Nota sobre la referencia de tempo}
El criterio original exigía que al menos 2 clips por género incluyeran una
referencia de tempo en BPM numérico. Sin embargo, al inspeccionar el corpus
se verificó que \emph{solo 1 de los 5\,521 clips} de MusicCaps contiene un BPM
numérico explícito; los \textit{captions} describen el tempo de forma
\emph{cualitativa} (p.\,ej.\ \textit{``fast tempo''}, \textit{``slow
groove''}), con 2\,497 clips que mencionan el tempo verbalmente.
Por tanto, el criterio se adapta de la siguiente forma: el filtro de
\textit{caption} exige una \emph{mención} de tempo (cualitativa o numérica),
y la columna \texttt{tempo\_bpm} del banco se completa con una
\emph{estimación} de BPM derivada del propio audio mediante
\texttt{librosa} (\textit{beat tracking}), marcada como estimada.
Esta decisión preserva la utilidad de la anotación de tempo sin fabricar
valores inexistentes en la fuente.

\subsubsection{Disponibilidad de candidatos por género}

Tras aplicar el filtro de \textit{caption} (longitud, género e instrumento)
sobre los 5\,521 clips, el número de candidatos válidos por género es amplio
en todos los casos, lo que garantiza poder muestrear 10 clips accesibles por
género (Tabla~\ref{tab:candidatos}).

\begin{table}[t]
  \caption{Candidatos válidos por género tras el filtro de \textit{caption}}
  \label{tab:candidatos}
  \centering
  \begin{tabular}{lr}
    \hline
    \textbf{Género} & \textbf{Candidatos} \\
    \hline
    Electrónica & 796 \\
    Rock        & 514 \\
    Clásica     & 208 \\
    Pop         & 179 \\
    Folk        & 175 \\
    Jazz        & 160 \\
    Hip-Hop     & 98 \\
    Latina      & 68 \\
    R\&B        & 63 \\
    Metal       & 43 \\
    \hline
  \end{tabular}
  \smallskip\\
  \footnotesize{De cada conjunto se muestrean 10 clips accesibles
  (semilla\,=\,42).}
\end{table}

\subsubsection{Procedimiento de descarga y preprocesamiento}

El guion \texttt{sample\_testset.py} implementa el siguiente pipeline
reproducible (semilla fija\,=\,42):

\begin{enumerate}
  \item Descarga el CSV oficial \texttt{musiccaps-public.csv}.
  \item Filtra los \textit{captions} por longitud ($\geq 20$ palabras) y
        presencia de género e instrumento.
  \item Estratifica los candidatos por género a partir del campo
        \texttt{aspect\_list}.
  \item Selecciona aleatoriamente (semilla\,=\,42) 10 clips por género; la
        verificación de disponibilidad es \emph{perezosa}: se prueban
        candidatos barajados con \texttt{yt-dlp -{}-simulate} hasta lograr 10
        accesibles por género.
  \item Descarga el tramo de 10\,s con
        \texttt{yt-dlp -{}-download-sections} y \texttt{ffmpeg}.
  \item Normaliza a 32\,kHz mono con
        \texttt{torchaudio.transforms.Resample}.
  \item Verifica la integridad: cada archivo \texttt{.wav} debe contener
        \emph{exactamente} 320\,000 muestras ($10\,\text{s} \times
        32\,\text{kHz}$).
  \item Estima el tempo (BPM) del audio con \texttt{librosa} y lo registra en
        la metadata.
\end{enumerate}

\subsubsection{Entregable}

El proceso produce un directorio \texttt{testset/} con 100 archivos nombrados
como \texttt{NNNN\_genero\_BPM.wav}, un archivo \texttt{testset\_metadata.csv}
con las columnas
\texttt{[id, ytid, caption, genre, tempo\_bpm, instruments]}, y el guion
reproducible \texttt{sample\_testset.py} incluido en el repositorio.
Cada clip queda garantizado a 32\,kHz mono con exactamente 320\,000 muestras,
condición necesaria para el cálculo consistente de las métricas FAD, CLAP,
KLD y KAD de la Sección~\ref{sec:benchmark}.

% ============================================================
%  Referencias sugeridas adicionales (agregar al .bib)
% ============================================================
%
% @inproceedings{bogdanov2019mtg,
%   title={The {MTG}-{J}amendo Dataset for Automatic Music Tagging},
%   author={Bogdanov, Dmitry and Won, Minz and Tovstogan, Philip and
%           Porter, Alastair and Serra, Xavier},
%   booktitle={ICML Machine Learning for Music Discovery Workshop}, year={2019}}
%
% (agostinelli2023musiclm ya está declarada en benchmark_section.tex)
